% Owner: Role A. Translated and adapted from the verified Vietnamese draft
% docs/report_draft_b_c.md (section "b. Introduction"). Every claim in that
% source file was fact-checked against src/data.py, data/raw/data.csv, and
% outputs/results.csv before translation -- see progress/A.md 2026-08-30
% logs for the audit trail. If you edit this file, keep it consistent with
% the source draft or update both.

\section{Introduction}
\label{sec:b}

\subsection{Brief introduction to decision trees}

A decision tree is a supervised learning model used for both classification
and regression tasks. It represents the decision process as a tree
structure: each \emph{internal node} tests one feature (e.g.\ ``Tuition
fees up to date $\le 0.5$?''), each \emph{branch} corresponds to the
outcome of that test, and each \emph{leaf node} assigns a predicted class.

Training a decision tree follows a \textbf{greedy, top-down,
recursive-partitioning} strategy: at each node, the algorithm evaluates
every feature and every possible split threshold, and picks the split that
reduces class \emph{impurity} the most, then repeats recursively on each
resulting subset. Two impurity measures are standard:

\begin{itemize}
  \item \textbf{Entropy / Information Gain}: used by the ID3 algorithm
    \citep{quinlan1986induction}, measuring the information gained by a
    split.
  \item \textbf{Gini impurity}: used by CART \citep{breiman1984classification},
    measuring the probability of misclassification if a label were drawn
    at random from the node's class distribution.
\end{itemize}

The scikit-learn library used in this project implements an optimized
variant of CART \citep{pedregosa2011scikit}; it defaults to Gini but
supports switching to entropy via the \texttt{criterion} parameter.

The \textbf{main strength} of a decision tree compared to ``black-box''
models (neural networks, complex ensembles) is \textbf{interpretability}:
every path from the root to a leaf can be read directly as an
\texttt{IF (condition 1) AND (condition 2) ... THEN (label)} rule, with no
separate explanation technique (SHAP, LIME) required. Once nominal
categories have been encoded, trees also handle the resulting numerical
and indicator features without feature scaling. Both properties are used
directly in Section~\ref{sec:c} below.

The \textbf{main weakness} is a tendency to \textbf{overfit}: an
unconstrained tree keeps splitting until almost every leaf is pure on the
training set, effectively memorizing training noise, which yields very
high training accuracy but a much lower test accuracy
\citep{mitchell1997machine, breiman1984classification}. Section~\ref{sec:d}
shows this pattern directly in the baseline model, which motivates the
pruning improvement in Section~\ref{sec:f}.

\subsection{Objective of the project}

The project applies this theory to a real, socially meaningful problem:
\textbf{predicting whether a student drops out, remains enrolled longer
than expected, or graduates on time}, using real data from 4{,}424
university students in Portugal. Concretely, the group:

\begin{enumerate}
  \item Builds and evaluates an unconstrained baseline decision tree,
    measuring performance with the full set of metrics appropriate for a
    3-class classification problem (Accuracy, Precision, Recall,
    F1-score, ROC-AUC, Confusion Matrix).
  \item Analyzes the baseline tree's structure: what the root split is
    and why the algorithm chose it, whether the tree shows signs of
    overfitting, and extracts concrete decision rules.
  \item Proposes and tests three improvement methods, each targeting a
    different weakness of the baseline:
  \begin{itemize}
    \item \textbf{Pruning} (cost-complexity pruning): trades tree
      complexity for generalization to address overfitting.
    \item \textbf{Class-imbalance handling}
      (\texttt{class\_weight='balanced'}, SMOTE): addresses the model's
      default bias toward the majority class (Graduate). At baseline,
      Enrolled recall (0.3836) is well below Dropout's already-reasonable
      recall (0.6796), so Enrolled is the main target of this improvement.
    \item \textbf{Feature selection for early warning}: addresses a
      practical limitation. The baseline's strongest features are only
      known \emph{after} a semester ends, so the model with the highest
      accuracy is not necessarily the model that can actually be deployed
      for early intervention.
  \end{itemize}
  \item Quantitatively compares all five model configurations and asks
    \textbf{which model suits which application goal}, in line with the
    assignment brief's emphasis on explaining how the tree is built, how
    it behaves, and why each improvement helps or does not, rather than
    only reporting a final accuracy score.
\end{enumerate}
